% target: sec/04-family.tex
\subsection{The Family Beyond the Corpus}
\label{sec:beyond}

\src{A}{S}The systematic search (Section~\ref{sec:search}) shows how much of the literature the 37 methods stand for. \src{S}{}The screen includes <<included>> records that propose a method of the kind this survey covers<<?unsure_to_fulltext>>, and <<unsure_to_fulltext>> more pass it as unsure<</unsure_to_fulltext>>. Set against the methods of the corpus by the periods of the evaluation audit, the corpus holds <<corpus_P1>> methods against <<inc_P1>> included records up to 2009, <<corpus_P2>> against <<inc_P2>> from 2010 to 2015, <<corpus_P3>> against <<inc_P3>> from 2016 to 2020, and <<corpus_P4>> against <<inc_P4>> since 2021 (counting <<year2027>> forthcoming records dated <<year_future>>). \src{A}{S}The corpus thus covers the early literature closely and the recent literature thinly: for each corpus method of the last period, the screen includes about <<ratio_P4>> records.

\src{A}{S}The codes of the screen, given from titles and abstracts, show what the included records propose. \src{S}{}New variants of SMOTE are the largest group in every period, <<newvar_P1_pct>>\% of them up to 2009 and <<newvar_P4_pct>>\% since 2021, and hybrids that embed SMOTE in another procedure a falling share (<<hybrid_P1_pct>>\% and <<hybrid_P4_pct>>\%). \src{A}{S}Generative oversamplers, which learn the minority distribution with a GAN, an autoencoder, a diffusion or a language model, rose from <<gen_P1_pct>>\% and <<gen_P2_pct>>\% of the included records in the first two periods to <<gen_P3_pct>>\% from 2016 to 2020 and <<gen_P4_pct>>\% since 2021 (<<gen_P4>> records). The corpus counts one learned generator in the family, GAMO (Section~\ref{sec:definition}), and places the deep generative models beyond its boundary (Section~\ref{sec:methodview}); Section~\ref{sec:beyondtables} describes that branch.

\src{S}{}Whether the findings of Sections~\ref{sec:bydecision} to~\ref{sec:methodview} hold beyond the corpus is tested on the coverage sample (Section~\ref{sec:search}), whose papers are coded in full text as the corpus was. Of its <<cov_drawn>> papers, <<cov_assessed>> could be assessed and <<cov_eligible>> are eligible (<<cov_core>> Core, <<cov_partial>> Partial, and <<cov_notcounted>> Not counted); Figure~\ref{fig:coverage} sets them against the corpus, and the online appendix (Section~S13) gives the counts. The sample shows a family of the same shape: most methods are Core, most work on two classes only, and about half of the papers state that the resampling stayed inside the training folds. \src{S}{}It differs from the corpus as follows. Its variants change fewer of the six decisions, <<cov_dec_mean>> on average against <<corpus_dec_mean>> in the corpus, with the largest gaps in seed selection and sample allocation. More of its methods generate in another way than by interpolation (<<cov_notcounted_pct>>\% against <<corpus_notcounted_pct>>\%), most of them published since 2016. Its evaluations are lighter: fewer of its papers repeat their partitions, test significance, or release code. On one point, the sample does better than the corpus: fewer of its papers choose their settings on test scores. The corpus thus describes the methods of the family fairly, while its evaluations are more elaborate than those of the wider literature.

\begin{figure}[tp]
\centering
\includegraphics[width=0.98\linewidth]{figures/fig_coverage.pdf}
\caption{\src{S}{}The corpus against the coverage sample of the systematic search (Section~\ref{sec:search}), coded alike: the share of papers with each property. Filled circles: the corpus. Open circles: the eligible papers of the sample, assessed in full text, with a Wilson 95\% interval. The decisions are shares of the Core and Partial methods. The online appendix (Section~S13) gives the counts.}
\Description{\src{S}{}A dot plot with one row for each of twenty properties in four groups (method, decisions changed, class category, evaluation). In each row, a filled circle marks the share of corpus papers with the property and an open circle, with a horizontal line for its 95\% interval, the share in the coverage sample. The two agree on most method and class properties; the corpus is higher on seed selection, sample allocation, repeated partitions, settings chosen on test scores, significance tests, and code release.}
\label{fig:coverage}
\end{figure}
